\documentclass[11pt]{article}
\usepackage[final]{acl}
\usepackage{times}
\usepackage{latexsym}
\usepackage[T1]{fontenc}
\usepackage[utf8]{inputenc}
\usepackage{microtype}
\usepackage{amsmath,amssymb}
\usepackage{booktabs}
\usepackage{graphicx}

\title{Always Late: What Happens to Language Features\\When You Compress a Model}
\author{Hemang Joshi \and Eashaan Thakur \and Shardul Joshi \and Nikhita Ravi\\
International Institute of Information Technology, Hyderabad}

\begin{document}
\maketitle

\begin{abstract}
We looked inside two Gemma models with a sparse autoencoder and asked what happens to language-specific features when the weights are quantized. We found three things that come apart under compression: whether a feature still fires the same way (activation survival), whether removing it still hurts the language (functional effect), and whether it still does the same thing to the output (effect identity). The standard activation cutoff misclassifies about a third of functionally intact features. The features matter for text prediction in all 27 languages we tested, but they carry almost none of the actual quantization damage, and they do not matter for reading comprehension. GPTQ roughly halves the amplification effect compared to round-to-nearest. The pattern holds on Gemma-2-9B but not on Llama-3.2-1B.
\end{abstract}

\section{The problem}

When you round a model's weights to fewer bits, its predictions get worse. People have noticed that this hits some languages harder than others, especially lower-resource ones. Meanwhile, sparse autoencoders (SAEs) have found features inside these models that fire specifically for individual languages. So a natural question: does quantization damage run \textit{through} the loss of these language features?

Previous work by Duan (2026) measured how many language-selective SAE features ``survive'' quantization by checking whether their activations still correlate with the original. We wanted to go further and ask whether a feature that changed its activation pattern still \textit{does the same thing} when you edit it.

\section{Setup}

\paragraph{Models.} We used Gemma-2-2B as the main model, Gemma-2-9B as a confirmation, and Llama-3.2-1B as a cross-check. Each model has a public SAE (Gemma Scope for both Gemma sizes; a TopK SAE for Llama). We read the hidden state at roughly the midpoint of each model.

\paragraph{Languages.} 27 languages from FLORES-200, split evenly across three resource tiers (high, mid, low). Feature selection uses the dev split; evaluation uses 300 devtest sentences per language.

\paragraph{Quantization.} We tested three quantizers, all simulated (weights rounded but stored as floats, no integer kernels):

\begin{itemize}
\item \textbf{Legacy RTN}: round-to-nearest with shared attention scales, from the paper we replicated.
\item \textbf{Matched RTN}: standard round-to-nearest, one scale per output channel. Same rounding grid as GPTQ but without calibration.
\item \textbf{GPTQ}: calibrated rounding using 8{,}192 Wikipedia tokens per language (221k total), which adjusts each weight to compensate for earlier rounding errors.
\end{itemize}

\noindent We tested bit widths from 8 down to 4, with INT16/INT12 as rounding-noise floors and INT0 (identity) as a sanity check.

\section{Methodology}

\subsection{Finding language features}

A feature $f$ fires on a token when its SAE activation $a_f(t)>0$. We compute a selectivity score: the ratio of a feature's highest per-language firing rate to its average across languages. A feature counts as language-selective if this score exceeds the 95th percentile of the same score under shuffled language labels (the cutoff was 1.83). This gave us 8{,}086 selective features on Gemma-2-2B and 5{,}218 on Gemma-2-9B.

For each language, we take its top $K{=}64$ most selective features as the ``language group.'' We also select activity-matched control features that fire at similar rates but are not language-specific.

\subsection{Key metrics}

\paragraph{Activation survival.} The standard measure from Duan (2026). For feature $f$ at bit width $b$, compute the Pearson correlation of activations across tokens:
\begin{equation}
R_1(f, b) = \mathrm{corr}_t\!\big(a_f^{\mathrm{FP}}(t),\; a_f^{b}(t)\big)
\end{equation}
A feature ``survives'' if $R_1 > 0.9$.

\paragraph{Net harm.} We zero-ablate each language's group (set those features' contributions to zero in the residual stream) and measure the increase in next-token loss. Net harm subtracts the matched-control effect:
\begin{equation}
E_\ell = \underbrace{\Delta\mathrm{NLL}(\text{group}_\ell)}_{\text{group harm}} - \underbrace{\overline{\Delta\mathrm{NLL}(\text{controls}_\ell)}}_{\text{control harm}}
\end{equation}
where $\Delta\mathrm{NLL}$ is the change in negative log-likelihood (nats/token) on language $\ell$'s own text. Positive $E$ means the language group matters more than random features of similar activity.

\paragraph{Retention ratio.} To see whether the dependence on the features survives quantization:
\begin{equation}
\rho(b) = \frac{E^{(b)}}{E^{(\mathrm{FP})}}
\end{equation}
If $\rho \approx 1$, the features matter about as much in the quantized model as in full precision. If $\rho > 1$, the ablation effect is amplified (the compressed model depends on them more, relatively speaking).

\paragraph{Effect fingerprint.} To ask whether a feature still \textit{does the same thing}, we switch it off and record the centered change in output logits. The cosine similarity between the full-precision and quantized fingerprints gives:
\begin{equation}
W_1(f, b) = \frac{\mathbf{d}_f^{\mathrm{FP}} \cdot \mathbf{d}_f^{b}}{\|\mathbf{d}_f^{\mathrm{FP}}\|\;\|\mathbf{d}_f^{b}\|}
\end{equation}
where $\mathbf{d}_f$ is the change in logits from ablating feature $f$. A cosine of 1 means the feature does exactly the same thing; we compare this against a ``numerical floor'' from INT16 rounding.

\paragraph{Rescue fraction.} We patch full-precision activations back into the quantized model and measure what share of the quantization loss gap is recovered:
\begin{equation}
\text{rescue}(S, b) = \frac{\mathrm{NLL}^{(b)} - \mathrm{NLL}^{(b, \text{patch } S)}}{\mathrm{NLL}^{(b)} - \mathrm{NLL}^{(\mathrm{FP})}}
\end{equation}
where $S$ is the set of features being restored.

\section{A correction we found}

We reproduced Duan's numbers with their released code (INT6 survival 51.15\% vs their 51.30\%), then traced the large gap with our pipeline to a missing beginning-of-sequence (BOS) token in their Gemma input protocol. Adding BOS dropped full-precision WikiText perplexity from 458.4 to 16.1 and raised INT6 survival from 51\% to 97.5\%. We adopted the corrected protocol. We also retired our own first pipeline after finding that it counted BOS in feature statistics.

\section{Results}

\subsection{The features matter, and the effect survives compression}

Zero-ablating a language's group produces a median net harm of 0.262 nats/token at full precision, positive in all 27 languages. The retention ratio $\rho$ stays near 1.0 through INT5 and jumps at INT4:

\begin{table}[h]
\centering\small
\begin{tabular}{lcccccc}
\toprule
& INT8 & INT7 & INT6 & INT5 & INT4 \\
\midrule
$\rho$ (2B) & 1.01 & 1.03 & 1.04 & 1.13 & 2.14 \\
$\rho$ (9B) & 1.00 & 1.05 & 1.01 & 1.21 & 2.48 \\
\bottomrule
\end{tabular}
\caption{Retention ratio across precisions. Rounding-only controls (INT16, INT12) give exactly 1.00.}
\end{table}

The INT4 amplification happens because the quantized model is already damaged, not because the features carry the damage. The rescue study confirms this: restoring the language group's activations at INT4 recovers only 2.6\% of the loss gap, while restoring the whole layer residual recovers 39.3\%.

\subsection{Activation survival overstates the damage}

At legacy INT4, only 13.2\% of a 1{,}487-feature sample pass the $R_1 > 0.9$ cutoff. But by output effect, 99.7\% are still individually identifiable ($W_1$ retrieval). The activation cutoff misclassifies about 31\% of functionally intact features at INT5.

The INT4 loss of similarity turns out to be generic. We ran a control where we ablated a \textit{different} feature's direction at the same positions and amplitudes. The control's median $W_1$ at INT4 was 0.828 vs the real feature's 0.842. The gap (0.014) is tiny and consistent at every precision. So the similarity loss reflects how the downstream computation drifts under quantization for any perturbation, not anything specific to the feature.

\subsection{GPTQ helps substantially}

Comparing the three quantizers on Gemma-2-9B at INT4:

\begin{table}[h]
\centering\small
\begin{tabular}{lccc}
\toprule
& Legacy & Matched & GPTQ \\
\midrule
Survival (\%) & 22.3 & 32.9 & 46.6 \\
Loss (nats/tok) & 0.410 & 0.287 & 0.088 \\
Amplification $\rho$ & 1.69 & 1.49 & 1.17 \\
\bottomrule
\end{tabular}
\caption{Three quantizers on Gemma-2-9B at INT4. GPTQ roughly halves the amplification and more than doubles survival.}
\end{table}

\noindent Matched RTN alone (better implementation, same method) accounts for about a third of GPTQ's improvement. At INT6, all three quantizers produce retention near 1.0 and survival above 95\%.

\subsection{The Belebele null}

This was our only confirmatory (pre-registered) test. Zeroing a language's group does not change Belebele reading-comprehension accuracy beyond matched controls: $-0.0001$ nats/item $[-0.0035, +0.0031]$. The groups are small (median 192 features) and the edit path itself works (random sets of 1{,}024+ features do lower accuracy). The groups just do not carry enough weight for this task.

\subsection{Gemma-2-9B confirmation}

We ran the full study on Gemma-2-9B with all three quantizers. Both tiers of predictions, frozen before any 9B number existed, were met: tier~1 (Study~G, the loss-based tests) and tier~2 (Phase~1, survival and loss ordering). GPTQ was below legacy RTN in 26 of 27 languages on the amplification ratio. The 9B results confirm the 2B pattern within the same model family and SAE.

\subsection{Llama-3.2-1B: not consistent}

On Llama with a different SAE architecture (TopK vs JumpReLU), different training data, and a different tokenizer, the group effect exists in 21 of 27 languages but is small (median $E = 0.048$ vs Gemma's 0.262), and the INT4 amplification is absent (median ratio 0.86 vs Gemma's 1.54). Pre-registered verdict: not consistent. But model family, SAE architecture, recipe, and tokenizer all change together, so this does not isolate what caused the difference.

\section{What it means}

The project's main contribution is methodological. Three things that sound like they should go together actually come apart under quantization:

\begin{enumerate}
\item \textbf{Activation survival} ($R_1 > 0.9$): whether the feature fires the same way.
\item \textbf{Functional effect} ($E$): whether removing it still hurts the language.
\item \textbf{Effect identity} ($W_1$): whether it still does the same thing to the output.
\end{enumerate}

The activation cutoff overstates functional damage. The INT4 effect-identity loss is generic across directions. And the amplification of the ablation effect does not mean the features carry the quantization damage; the rescue study directly measures this at 2.6\%.

On the practical side, GPTQ's improvement is large and consistent across model sizes: it roughly halves the INT4 amplification and more than doubles survival. The matched RTN comparison shows that about a third of this comes from a better rounding implementation alone, without calibration data.

The Belebele null and the Llama inconsistency are honest limits. One confirmatory test produced a null, and the cross-family check failed (though it is heavily confounded). The findings hold within Gemma~2 with Gemma Scope at one layer per model size. Whether they generalize to other architectures, layers, or SAE families is open.

\end{document}
